% !TEX program = xelatex
\documentclass[a4paper]{article}
\usepackage{amsthm}
\usepackage{amssymb}
\usepackage{bm}
\usepackage{mathtools}
\usepackage[x11names]{xcolor}
\usepackage{xparse}
\usepackage{fontspec}
\usepackage{unicode-math}
\setromanfont{DovesType-Regular.otf}
\setsansfont{Andika}
\setmathfont{Asana Math}[Scale=1]

% \usepackage{pstricks}
\usepackage{varwidth}
\usepackage{siunitx}
\usepackage{graphicx}
\usepackage[margin=1cm]{geometry}
\usepackage[most]{tcolorbox}
\usepackage{pgfplots}
\pgfplotsset{compat=newest}
\tcbuselibrary{skins,xparse,poster,breakable}
% \usetikzlibrary{fadings}
\usetikzlibrary{calc, plotmarks, shapes, shapes.geometric, positioning, angles, intersections, quotes, through, patterns, turtle, arrows.meta}
\usetikzlibrary{decorations.markings,backgrounds}
% \usepackage{etoolbox}
% \usepackage{tkz-euclide}
% \usepackage{xlop}
% \newcommand\hole[2]{#1}  % for use with xlop
\pagenumbering{gobble}
%%%%%%%%%%%%%%%%%%%%%%%%%%%%%%%%%%%%%%%%%%%%%%%%%%%%%%%%%
\newcommand\markangle[9]{% origin X Y radius radiusmark mark colour opacity
%  % fill red circle offset-from-centre
  \begin{scope}
    \path[clip] (#1) -- (#2) -- (#3);
    \fill[color=#7,fill opacity=#8,draw=black,name path global=pcircle]  % global declaration required otherwise pcircle is not seen by the `named intersections=' lines below.
    (#1) circle (#4);
  \end{scope}
  % middle calculation
  \path[name path=line one] (#1) -- (#2);
  \path[name path=line two] (#1) -- (#3);
  \path[%
  name intersections={of=line one and pcircle, by={inter one}},
  name intersections={of=line two and pcircle, by={inter two}}
  ] (inter one) -- (inter two) coordinate[pos=#9] (place);
  % put mark
  \node at ($(#1)!#5!(place)$) {\scriptsize{#6}};
}

\newcommand\tcircle[6]{% centre coord (x,y), radius, points, radpoint, colour, edge
  \coordinate (O) at (#1,#2); % centre of the circle
  \def\radius{#3}          % radius of the circle
  \def\npts{#4}            % number of the points
  \def\radpt{#5}           % radius of the points
  \colorlet{ptcolour}{#6}  % colour of the points
  % \draw (O) circle (\radius);
  \foreach \numpoint in {1,...,\npts}{
    \fill[ptcolour] (O) ++ (360/\npts*\numpoint:\radius) coordinate (C\numpoint) circle(\radpt);
  }
}

% \newcommand{\condSoln}[2]{\ifcsdef{r@#1}{#2}{}}

% \newcommand\fadingtext[3][]{%
%    \begin{tikzfadingfrompicture}[name=fading letter]
%      \node[text=transparent!0,inner xsep=0pt,outer xsep=0pt,#1] {#3};
%    \end{tikzfadingfrompicture}%
%    \begin{tikzpicture}[baseline=(textnode.base)]
%      \node[inner sep=0pt,outer sep=0pt,#1](textnode){\phantom{#3}};
%      \shade[path fading=fading letter,#2,fit fading=false]
%      (textnode.south west) rectangle (textnode.north east);%
%    \end{tikzpicture}%
% }

\definecolor{JISpurple}{RGB}{89,72,122}
\definecolor{JISivory}{RGB}{241,234,221}
\definecolor{JIStaupe}{RGB}{183,156,154}
\definecolor{PaleGreen}{RGB}{240,255,240} % 'Honeydew'

\AddToHook{shipout/background}{%
    \put (0in,-\paperheight){\includegraphics[width=\paperwidth,height=\paperheight]{images/R10V5-R.png}}%
}

\newcommand\numberthis{\addtocounter{equation}{1}\tag{\theequation}}

\newtcolorbox{MyOuterBox}{%
  enhanced,
  breakable,
  frame style=JISpurple,
  colback=JISivory,
  colframe=JISpurple,
  title={\includegraphics[width=0.9cm,height=0.9cm]{images/JIS Final Logo FA-02.png}\raisebox{3mm}{\Large{Maths Challenge}\hspace{26em} \Large{\bfseries\sffamily 13}}},
}

\newtcolorbox{MyInnerBox}[2][]{enhanced,%empty,
coltitle=JISpurple,colback=white,
breakable,
fonttitle=\bfseries\sffamily,
attach boxed title to top left={yshift=-1.5mm},
boxed title style={empty, size=small, top=1mm, bottom=0pt},
varwidth boxed title=0.5\linewidth,
frame code={
  \path (title.east|-frame.north) coordinate (aux);
\path[draw=JISpurple, line width=0.5mm, rounded corners,fill=white]
(frame.west) |- ([xshift=-2.5mm]title.north east) to[out=0, in=180] ([xshift=7.5mm]aux)-|(frame.east)|-(frame.south)-|cycle;
},
title={#2},#1}

\newtcolorbox{MyInnerSplitBox}[2][]{enhanced,%empty,
bicolor,sidebyside,sidebyside align=top seam,
righthand width=8.5cm,colbacklower=white,
sidebyside gap=5mm,
breakable,
coltitle=JISpurple,colback=white,
fonttitle=\bfseries\sffamily,
attach boxed title to top left={yshift=-1.5mm},
boxed title style={empty, size=small, top=1mm, bottom=0pt},
varwidth boxed title=0.5\linewidth,
frame code={
  \path (title.east|-frame.north) coordinate (aux);
\path[draw=JISpurple, line width=0.5mm, rounded corners,fill=white]
(frame.west) |- ([xshift=-2.5mm]title.north east) to[out=0, in=180] ([xshift=7.5mm]aux)-|(frame.east)|-(frame.south)-|cycle;
},
title={#2},#1}


\newtcolorbox{MySolutionBox}[1][]{%
  enhanced,
  breakable,
  frame style=JISpurple,
  colback=PaleGreen, colframe=green,
  title={\Large Solution},
  drop fuzzy shadow,
  halign=left,
  #1
}

%%%%%%%%%%%%%%%%%%%%%%%%%%%%%%%%%%%%%%%%%%%%%%%%%%
\newtoggle{SOLUTION}
%%% Uncomment the appropriate line below to show solutions %%%
% \toggletrue{SOLUTION}
\togglefalse{SOLUTION}
%%%%%%%%%%%%%%%%%%%%%%%%%%%%%%%%%%%%%%%%%%%%%%%%%


%%%%%%%%%%%%%%%%%%%%%%%%%%%%%%%%%%%%%%%%%%%%%%%%%%
%%%%%%            DOCUMENT BEGINS           %%%%%%
%%%%%%%%%%%%%%%%%%%%%%%%%%%%%%%%%%%%%%%%%%%%%%%%%%
\begin{document}


  \begin{MyOuterBox}
    \iftoggle{SOLUTION}{Here are the full, or partial solutions.
    }{
      Welcome to this week's Maths Challenge!\\
      It's fine to have a go at both questions!\\
      Solutions must be explained in detail, responses with just the answer may be ignored.\\
      Drop your solution in the box in the staffroom by Tuesday.
    }
       \begin{MyInnerBox}{Year 8 and below}
      Thirty-one people took an exam and each person achieved a different, whole number score, from \(70\) to \(100\). The average of the scores is calculated. Then one person's score is removed, but when the average is recalculated, it has not changed, the new mean is the same as the mean of the original \(31\) marks. Which score was taken out?
      \iftoggle{SOLUTION}{%conditional output begin
      \begin{MySolutionBox}
        It might be easier to start with an example containing easier numbers. Suppose we have five examinees and their scores are 1, 2, 3, 4, 5.\par
        The sum of the scores is 15 and the mean is \(\frac{15}{5}=3\).\par
        If we take out the middle score, 3, the new sum of scores is 12 and the new mean is \(\frac{13}{4}=3\), the mean did not change when we removed the middle value.\par
        Notice also that if we find the sum by adding in pairs, \(1+5 = 6\), then \(2+4=6\) until the middle value is left, the total is the same each time, there is a kind of symmetry around the middle value.\par
        Now if we start with the problem given, we can add up the thirty-one scores in pairs: \(70+100=170\), then \(71+99=170\), then \(72+98=170\) and so on until we reach \(84+86=170\), that's fifteen pairs, leaving \(85\) as the middle number.\par
        The average of the thirty-one scores:
        \begin{align*}
          \frac{(170\times 15)+85}{31} &= 85\\
          (170\times 15) + 85 &= 31 \times 85\\
          (170\times 15) + 85 &= 30 \times 85 + 85\\
          170\times 15 &= 30\times 85
          \shortintertext{The average of the thirty scores after removing the middle value, \(85\):}
          \frac{170\times 15}{30} &= 85
        \end{align*}
      The averages are the same.\par
        \(85\) was the score removed.
      \end{MySolutionBox}
    }{}%conditional output end
    \end{MyInnerBox}


    \vspace{0.4cm}
          \begin{MyInnerBox}{Year 9 and above}
        \begin{minipage}[t]{0.5\linewidth}
          \strut\vspace*{-\baselineskip}\newline
        The diagram shows a trapezium. The line of length \(\SI{3}{\cm}\) is perpendicular to the two parallel lines of the trapezium. The purple area is \(\SI{6}{\square\cm}\) greater than the green area. Find the length \(x\).\par
        \end{minipage}%
        \begin{minipage}[t]{0.5\linewidth}
          \strut\vspace*{-\baselineskip}\newline
        \begin{tikzpicture}[scale=0.80]
          \draw (0,0) coordinate (A) -- (9,0) coordinate (B) -- (9,6) coordinate (C) node[right,midway] {\(3\)} -- (3,6) coordinate (D)  -- cycle node[left,midway] {\(x\)};
          \draw[name path=AC] (A) -- (C);
          \draw[name path=BD] (B) -- (D);
          \path[name intersections={of=AC and BD,by=M}];
          \filldraw[black,fill=Green2, line join=bevel] (C) -- (D) -- (M) -- cycle;
          \filldraw[black,fill=DarkOrchid1, line join=bevel] (A) -- (B) -- (M) -- cycle;
          \iftoggle{SOLUTION}{
            \node[below left] at (A) {\(A\)};
            \node[below right] at (B) {\(B\)};
            \node[above right] at (C) {\(C\)};
            \node[above left] at (D) {\(D\)};
            \node[left=0.2] at (M) {\(M\)};
            \node[below] at (4.5,0) {\(a\)};
            \node[above] at (6,6) {\(b\)};
            \draw[dashed,Stealth-Stealth] let \p1=(M) in (9.7,0) -- (9.7,\y1) node[midway,fill=white] {\(h\)};
            \draw[dashed,Stealth-Stealth] let \p1=(M) in (9.7,\y1) -- (9.7,6) node[midway,fill=white] {\(3-h\)};
            % \node[] (H) at (M-|B) {*}; % can use this to get y-coord of M
            \draw[dashed] (3,0) node[below] {\(N\)} -- (D);
            \markangle{A}{C}{B}{6mm}{0mm}{}{yellow}{1}{1}
            \markangle{C}{A}{D}{6mm}{0mm}{}{yellow}{1}{1}
            \markangle{B}{A}{D}{6mm}{0mm}{}{Tomato1}{1}{1}
            \markangle{D}{B}{C}{6mm}{0mm}{}{Tomato1}{1}{1}
          }{}
        \end{tikzpicture}
      \end{minipage}\par%
      \iftoggle{SOLUTION}{%conditional output begin
      \begin{MySolutionBox}
        Surprisingly, we don't need to know much about trapeziums, other than that \(AB\parallel DC\), to solve this interesting geometry puzzle. It's more important to remember similar triangles, triangle area, and good old Pythagoras. To help us proceed, we add some labels, the vertical height \(h\) of the point \(M\) where the two diagonals intersect, and we make a right-angled triangle \(\bigtriangleup ADN\).\par
        By Pythagoras:
        \begin{align}
          x^{2} &= 3^{2} + (a-b)^{2}
          \shortintertext{Area \(P\) of the purple triangle, \(\bigtriangleup ABM\):}
          P &= \frac{1}{2} ah
          \shortintertext{Area \(G\) of the green triangle, \(\bigtriangleup CDM\):}
          G &= \frac{1}{2} b(3-h)
          \shortintertext{but we know \(P = G + 6\), use Equations 2 \& 3:}
          \frac{1}{2}ah &= 6 + \frac{1}{2}b(3-h)\notag\\
          ah &= 12 + 3b - bh\notag\\
          (a+b)h &= 12 + 3b
          \intertext{Now \(\angle CAB = \angle ACD\) and \(\angle ABD = \angle CDB\) so the purple triangle\(\bigtriangleup ABM\) and the green triangle \(\bigtriangleup CDM\) are similar. So equivalent distances in each triangle differ by the same proportion. For example DM is to BM as \(b\) is to \(a\). Also CM is to AM as \((3-h)\) is to \(h\). Therefore,}
          a:b &:: h:(3-h) \notag
          \shortintertext{that is,}
          \frac{a}{b} &= \frac{h}{3-h} \notag\\
          3a - ah &= bh \notag\\
          (a+b)h &= 3a
          \intertext{Now we set Equation 4 equal to Equation 5, eliminating \(h\):}
          3a &= 12 + 3b \notag\\
          a - b &= 4
          \intertext{now we can use Equation 6 to substitute \(4\) for \(a-b\) in Equation 1.}
          x^{2} &= 9 + 4^{2} \notag\\
          x^{2} &= 25 \notag\\
          x &= \pm 5 \notag
          \shortintertext{but distance \(x\) must be positive, so}
          x &= 5 \notag
        \end{align}
      \end{MySolutionBox}
      }{}%conditional output end
    \end{MyInnerBox}


  \end{MyOuterBox}

%%%%%%%%%%%%%%%%%%%%%%%%%%%%%%%%%%%%%%%%%%%%%%%%%%
\end{document}



